% Part: second-order-logic
% Chapter: syntax-and-semantics

\documentclass[../../../include/open-logic-chapter]{subfiles}

\begin{document}

\olchapter{sol}{syn}{వాక్యనిర్మాణం, అర్థవిచారం}

\begin{editorial}
ఇందులో ఇప్పటివరకు ద్వితీయ-స్థాయి తర్కపు ప్రాథమిక వాక్యనిర్మాణాన్నీ,
అర్థవిచారాన్నీ వివరించాం. అధ్యాయంగా ఇది ఇంకా చిన్నదే. ద్వితీయ-స్థాయి
తర్కపు \tetoken{వ్యుత్పత్తి}{!!{derivation}} వ్యవస్థల గురించి చెప్పాలంటే, ద్వితీయ-స్థాయి చరాల
ప్రతిస్థాపనను వివరించాలి. అందులో కొన్ని సూక్ష్మమైన సమస్యలున్నాయి.
\end{editorial}

\olimport{introduction}

\olimport{terms-formulas}

\olimport{satisfaction}

\olimport{semantic-notions}

\olimport{expressive-power}

\olimport{inf-count}

\OLEndChapterHook

\end{document}
